\documentclass{article}
\usepackage{graphicx} % Required for inserting images
\usepackage{amssymb}

\title{Fundamentals Chapter 3}
\author{Siddharth K}
\date{October 2024}

\newcommand{\N}{\mathbb{N}}
\newcommand{\Z}{\mathbb{Z}}
\newcommand{\R}{\mathbb{R}}
\newcommand{\C}{\mathbb{C}}

\begin{document}

\maketitle

\textbf{1.} Consider the group $G = (\Z/p\Z)^\times$ under multiplication where $p$ is prime. The order of $G$ is $p - 1$. By Lagrange's theorem, the order of any subgroup of $G$ divides $|G|$. Now, consider $a \in G$ and the subgroup $\langle a \rangle \le G$ and let $|\langle a\rangle| = k$. We know that $a^k \equiv 1 \pmod{p}$, so
$$a^{p - 1} \equiv (a^k)^{\frac{p - 1}{k}} \equiv 1 \pmod{p}$$
as desired. Multiplying by $a$ gives us $a^p \equiv a \pmod{p}$. \\

\textbf{2.} There is no group $G$ and subgroup $H \le G$ where $(aH)(bH) = abH$ for all $a, b \in G$. To see this, let $b = a^{-1}$. Then $(aH)(a^{-1}H) = \{ah_1a^{-1}h_2: h_1, h_2 \in H\}$. By the condition, this is equal to $H$, which means that $ah_1a^{-1} \in H$, which implies that $H$ is normal. \\

\textbf{3.} Since $G$ is abelian, $ab = ba$ for all $a, b \in G$. So
$$aH = \{ah: h \in H\} = \{ha: h \in H\} = Ha,$$
which means $H$ is normal. \\

\textbf{5.} Let $K = \ker(\phi)$. Then for all $k \in K$, $\phi(k) = e_G$. Also, since $\phi$ is a homomorphism,
$$\phi(ak) = \phi(a)\phi(k) = \phi(a)e_G = \phi(a),$$
$$\phi(ka) = \phi(k)\phi(a) = e_G\phi(a) = \phi(a)$$
for all $a \in H$. So $\phi(ak) = \phi(ka)$. This implies that, for all $k_1 \in K$, there exists $k_2 \in K$ such that $ak_1 = k_2a$. Therefore, $aK = Ka$ and $K$ is a normal subgroup of $H$. \\

\textbf{7.} From problem 5, if there's a group $H$ and a homomorphism $\phi: G\to H$ then $\ker(\phi)$ is normal. \\

Assume $K \le G$ is normal. Now, define a map $\phi : G \to G/K$ by $\phi(g) = gK$ for all $g \in G$. This is a homomorphism because
$$\phi(g_1g_2) = (g_1g_2)K = (g_1K)(g_2K) = \phi(g_1)\phi(g_2)$$
for all $g_1, g_2 \in G$. \\

The kernel of $\phi$ is
$$\ker(\phi) = \{g \in G: \phi(g) = K\} = \{g \in G: gK = K\},$$
which simplifies to
$$\ker(\phi) = \{g \in G: g \in K\} = K,$$
so $\ker(\phi) = K$. \\

\textbf{9a.} We need to show that $g_1 \cdot (g_2 \cdot x) = (g_1g_2) \cdot x$ and $e \cdot x = x$ for all $g_1, g_2 \in G$ and $x \in G$. Looking at $g_1 \cdot (g_2 \cdot x)$, we get
$$g_1 \cdot (g_2 \cdot x) = g_1 \cdot (g_2xg_2^{-1})$$
$$= g_1(g_2xg_2^{-1})g_1^{-1}$$
$$= g_1g_2xg_2^{-1}g_1^{-1}.$$
Also, $e\cdot x = exe^{-1} = exe = x$, so $\cdot$ defines a group action. \\

\textbf{9b.} Since $C(x) \subset G$, we need to show that $e \in C(x)$ and $g^{-1} \in C(x)$ for all $g \in C(x)$. Since $ge = eg$ for all $g \in G$, $e \in C(x)$. \\

To show that $g^{-1} \in C(x)$, we need to prove that $g^{-1}x = xa^{-1}$. Since $a \in C(x)$, we know that $ax = xa$, so
$$e = (xa)(xa)^{-1}$$
$$= (ax)(xa)^{-1}$$
$$= axa^{-1}x^{-1}.$$
Multiplying both sides by $a^{-1}$ on the left and then $x$ on the right, we get
$$a^{-1} = xa^{-1}x^{-1}$$
$$a^{-1}x = xa^{-1},$$
which means $a^{-1} \in C(x)$. \\

\textbf{10a.} Consider the set:
$$aZ(G) = a\{x \in G: gx = xg\}$$
$$= \{ax \in G: gx = xg\}$$
$$= \{xa \in G: gx = xg\}$$
$$= Z(G)a,$$
so $Z(G)$ is normal. \\

\textbf{10b.} If $G/Z(G)$ is cyclic then there exists $gZ(G) \in G/Z(G)$ such that every coset in $G/Z(G)$ can be written as $(gZ(G))^k$ for some $k \in \Z$. So, for all $h \in G$,
$$hZ(G) = (gZ(G))^k = g^kZ(G)$$
for some $k \in \Z$. This means $h$ and $g^k$ are in the same coset of $Z(G)$ in $G$. This implies that $h(g^k)^{-1} \in Z(G)$ or $h = g^kx$ for some $x \in Z(G)$. \\

Now, let $a, b \in G$ be
$$a = g^{k_1}x_1, \quad b = g^{k_2}x_2.$$
Then
$$ab = g^{k_1}x_1g^{k_2}x_2$$
$$= g^{k_1}g^{k_2}x_1x_2$$
$$= g^{k_1 + k_2}x_1x_2$$
$$ba = g^{k_2}x_2g^{k_1}x_1$$
$$= g^{k_2}g^{k_1}x_2x_1$$
$$= g^{k_1 + k_2}x_1x_2$$
so $ab = ba$ for all $a, b \in G$ and $G$ is abelian. \\

\textbf{11.} To show that $|aH| = |bH|$ for any $a, b \in G$, we need to establish a bijection from $aH$ to $bH$. We can map each $ah \in aH$ to $ba^{-1}ah = bh \in bH$. Since multiplication by $ba^{-1}$ is unique, this is a bijection. \\

Now, we need to show that if $aH \neq bH$ then $aH \cap bH = \varnothing$. Assume that $aH \cap bH \neq \varnothing$. Then there exists $x \in G$ such that $x \in aH$ and $x \in bH$. This implies that $x = ah_1 = bh_2$, so $b^{-1}a = h_2h_1^{-1}$, so $b^{-1}a \in H$. This implies that $a \in bH$ and $aH = bH$, which is a contradiction. So $aH \cap bH = \varnothing$.

\end{document}